\chapter{Υλοποίηση}
\label{chap:implementation}

Το κεφάλαιο αυτό περιγράφει την υλοποίηση του συστήματος: τη δομή του πακέτου, τον μηχανισμό επαλήθευσης, τις συντεταγμένες και τους παραγόμενους πίνακες, τις εγγενείς μηχανές αναζήτησης και τη δοκιμαστική σουίτα. Έμφαση δίνεται στα σημεία όπου οι σχεδιαστικές αποφάσεις του Κεφαλαίου~\ref{chap:design} αποκτούν συγκεκριμένη μορφή σε κώδικα και σε παραγόμενα τεκμήρια.

\section{Δομή πακέτου}
\label{sec:implementation-package}
Το πακέτο Python ονομάζεται \code{rubik_optimal}. Τα βασικά αρχεία του είναι τα εξής:

\begin{itemize}
\item \code{cube.py}: αναπαράσταση κατάστασης στο μοντέλο κυβιδίων (\eng{cubies}), κινήσεις και έλεγχος φυσικής εγκυρότητας,
\item \code{moves.py}: συντακτική ανάλυση κινήσεων και αντίστροφες ακολουθίες,
\item \code{search/bfs.py}: πλήρης αναζήτηση κατά πλάτος (\eng{breadth-first search}, BFS) μικρού βάθους,
\item \code{search/ida_star.py}: επαναληπτικά βαθυνόμενη A* (\eng{Iterative-Deepening A*}, IDA*) με παραδεκτό κάτω φράγμα,
\item \code{native/corner_pdb/corner_pdb.cpp}: εγγενής γεννήτρια του γωνιακού πίνακα προτύπων (\eng{pattern database}, PDB) του 3×3×3,
\item \code{native/edge_pdb/edge_pdb.cpp}: εγγενής γεννήτρια των εξακμικών πινάκων προτύπων του 3×3×3,
\item \code{native/optimal_solver/optimal_solver.cpp}: εγγενής IDA* πλήρους κύβου, η οποία επιστρέφει ακριβείς (\code{exact}) λύσεις 3×3×3 όταν ολοκληρώνεται η αναζήτηση,
\item \code{native/h48_backend/h48_backend.c}: εγγενές περίβλημα πάνω στο ενσωματωμένο υποσύστημα H48 του \code{nissy-core} (άδεια GPL),
\item \code{tables/corner_pdb.py}: ανάγνωση του δυαδικού γωνιακού πίνακα με απεικόνιση μνήμης (\eng{memory mapping}),
\item \code{tables/edge_pdb.py}: αντίστοιχη ανάγνωση των δυαδικών πινάκων προτύπων ακμών,
\item \code{tables/h48.py}: μεταγλώττιση του υποσυστήματος H48, παραγωγή πινάκων, μεταδεδομένων και αθροισμάτων ελέγχου (\eng{checksums}),
\item \code{solvers/}: τυποποιημένα αποτελέσματα των επιλυτών (\eng{solvers}),
\item \code{solvers/h48_native.py}: υποσύστημα H48 με αποτελέσματα αποκλειστικά \code{exact} ή \code{timeout} και άμεση μετατροπή κατάστασης μεταξύ μορφής κυβιδίων και μορφής ψηφίδων (\eng{facelet}),
\item \code{benchmark.py}: παραγωγή αποτελεσμάτων, συνόψεων, πινάκων και δεδομένων σχημάτων,
\item \code{cli.py}: ενιαίο περιβάλλον γραμμής εντολών.
\end{itemize}

Η δομή αυτή κρατά χωριστές τρεις ευθύνες που συχνά συγχέονται σε μικρές υλοποιήσεις επιλυτών: το μοντέλο του κύβου, τις μηχανές αναζήτησης και την παραγωγή επιστημονικών αποτελεσμάτων. Το αρχείο του μοντέλου κυβιδίων δεν γνωρίζει τίποτα για μετρήσεις ή LaTeX. Οι επιλυτές δεν γράφουν απευθείας πίνακες του κειμένου. Η μονάδα μετρήσεων δεν επαναϋλοποιεί κινήσεις. Η διάκριση αυτή επιτρέπει να ελέγχεται κάθε επίπεδο με διαφορετικές δοκιμές.

Οι φάκελοι \code{coordinates/} και \code{tables/} αποτελούν το κεντρικό υπόστρωμα πινάκων της υλοποίησης. Οι κωδικοποιητές συντεταγμένων (\eng{coordinate encoders}) μετατρέπουν καταστάσεις κυβιδίων σε ακέραιες τιμές. Οι γεννήτριες πινάκων μετατρέπουν τις συντεταγμένες αυτές σε πίνακες κινήσεων (\eng{move tables}) και πίνακες αποκοπής (\eng{pruning tables}). Οι επιλυτές χρησιμοποιούν τις παραγόμενες πληροφορίες μόνο μέσω συναρτήσεων ευρετικής, ώστε να μη χρειάζεται κάθε επιλυτής να γνωρίζει λεπτομέρειες σειριοποίησης.

Το \code{results.py} παρέχει κοινές βοηθητικές συναρτήσεις ανάγνωσης και εγγραφής JSON και JSONL. Η ύπαρξη κοινής στρώσης για τα αρχεία αποτελεσμάτων μειώνει τον κίνδυνο να παραχθούν διαφορετικά σχήματα δεδομένων από διαφορετικά προγράμματα. Αυτό είναι πρακτικά σημαντικό, επειδή κάθε αριθμός του κειμένου πρέπει να είναι ανιχνεύσιμος σε συγκεκριμένο αποθηκευμένο αρχείο.

\section{Επαλήθευση λύσης}
\label{sec:implementation-verifier}
Ο μηχανισμός επαλήθευσης (\eng{verifier}) δεν εμπιστεύεται την εσωτερική δήλωση κανενός επιλυτή. Ξεκινά από την αρχική κατάσταση, εφαρμόζει τη δηλωμένη ακολουθία κινήσεων, μετρά το μήκος στη μετρική μισών στροφών (\eng{half-turn metric}, HTM) και ελέγχει αν η τελική κατάσταση είναι λυμένη. Κάθε γραμμή αποτελεσμάτων κρατά χωριστά τα πεδία \code{status} και \code{verified}.

Ο μηχανισμός αυτός έχει δύο διακριτούς ρόλους. Πρώτον, προστατεύει από σφάλματα επιλυτών: ένας επιλυτής μπορεί να επιστρέψει συντακτικά έγκυρη ακολουθία που δεν λύνει την κατάσταση, είτε λόγω σφάλματος αναζήτησης είτε λόγω ασυμφωνίας σύμβασης. Δεύτερον, προστατεύει από σφάλματα καταγραφής: μια γραμμή αποτελεσμάτων με \code{verified=true} πρέπει να μπορεί να επαληθευτεί ξανά αποκλειστικά από τα αποθηκευμένα πεδία της. Αυτό σημαίνει ότι το αποτέλεσμα δεν εξαρτάται από προσωρινή μνήμη Python ή από αντικείμενο που χάθηκε μετά την εκτέλεση.

Η επαλήθευση δεν αποδεικνύει βελτιστότητα. Αποδεικνύει ότι η ακολουθία λύνει την κατάσταση και ότι το δηλωμένο μήκος αντιστοιχεί στο πλήθος των αναγνωρισμένων κινήσεων. Για την ετικέτα \code{exact} απαιτείται επιπλέον απόδειξη από την ίδια την αναζήτηση, ενώ για την ετικέτα \code{non_exact} η επαλήθευση αρκεί για να τεκμηριώσει πρακτική λύση, όχι όμως ελάχιστο μήκος. Το διπλό αυτό επίπεδο είναι κεντρικό για την ακεραιότητα των ισχυρισμών της εργασίας και αναλύεται συνολικά στο Κεφάλαιο~\ref{chap:validation}.

\section{Ετικέτες κατάστασης}
\label{sec:implementation-status}
Οι ετικέτες κατάστασης που χρησιμοποιούνται είναι \code{exact} (αποδεδειγμένα βέλτιστη λύση), \code{non_exact} (επαληθευμένη αλλά όχι αποδεδειγμένα βέλτιστη λύση), \code{lower_bound} (αποδεδειγμένο κάτω φράγμα χωρίς λύση), \code{timeout} (διακοπή λόγω ορίων) και \code{not_applicable} (μη υποστηριζόμενη περίπτωση). Η ετικέτα \code{exact} σημαίνει ότι ολοκληρώθηκε εξαντλητική ή παραδεκτή αναζήτηση για τη συγκεκριμένη περίπτωση. Η ετικέτα \code{non_exact} σημαίνει ότι η λύση επαληθεύτηκε, αλλά δεν υπάρχει απόδειξη βελτιστότητας.

Η ετικέτα \code{lower_bound} χρησιμοποιείται όταν η διαδικασία δεν έχει λύση, αλλά έχει υπολογίσει παραδεκτό κάτω φράγμα. Η τιμή αυτή είναι χρήσιμη επειδή αποκλείει μικρότερες αποστάσεις, δεν είναι όμως η ίδια η απόσταση. Η ετικέτα \code{timeout} δηλώνει ότι η διαδικασία σταμάτησε από όριο χρόνου, κόμβων ή βάθους. Η ετικέτα \code{not_applicable} προορίζεται για επιλυτή που δεν υποστηρίζει συγκεκριμένη περίπτωση ή για ελλείπουσα προαιρετική εξάρτηση, ενώ η ετικέτα \code{failed} υπάρχει στο σχήμα για πραγματική αποτυχία, όπως μη επαληθευμένη ακολουθία. Το πλήρες σχήμα των αρχείων αποτελεσμάτων τεκμηριώνεται στο Παράρτημα~\ref{app:schema}.

Στα αποτελέσματα της εργασίας οι απαιτούμενες εγγενείς υλοποιήσεις δεν επιστρέφουν \code{not_applicable} ως βασική συμπεριφορά· η ετικέτα αυτή εμφανίζεται μόνο όταν λείπει προαιρετική εξάρτηση. Οι εξαντλήσεις χρονικού ορίου παραμένουν τεκμηριωμένος περιορισμός των βαθύτερων αναζητήσεων και αναλύονται στο Κεφάλαιο~\ref{chap:experiments}.

\section{Συντεταγμένες και πίνακες}
\label{sec:implementation-tables}
Η υλοποίηση περιλαμβάνει συντεταγμένες (\eng{coordinates}) προσανατολισμού γωνιών (CO), προσανατολισμού ακμών (EO), μετάθεσης γωνιών, μετάθεσης ακμών και φέτας UD (\eng{UD-slice}). Οι προβολές CO, EO και UD-slice χρησιμοποιούνται για τη φάση 1 του αλγορίθμου δύο φάσεων και για το γενικό παραδεκτό κάτω φράγμα της IDA*. Προστέθηκαν επίσης τρεις φασικές προβολές για τη φάση 2 του Kociemba: μετάθεση γωνιών, μετάθεση των οκτώ ακμών U/D και μετάθεση των τεσσάρων ακμών της φέτας. Οι φασικές αυτές προβολές παράγονται μόνο με το περιορισμένο σύνολο δέκα κινήσεων της φάσης 2, \(U,U',U2,D,D',D2,R2,L2,F2,B2\). Οι αντίστοιχοι πίνακες είναι επομένως μικρότεροι από έναν πλήρη πίνακα προτύπων του 3×3×3, αλλά ισχυρότεροι από την απλή μέτρηση λανθασμένων κυβιδίων.

Το πρόγραμμα \code{scripts/generate_tables.py} παράγει πίνακες κινήσεων και πίνακες αποκοπής για όλες τις παραπάνω προβολές, αποθηκεύοντας μέγεθος, χρόνο παραγωγής, άθροισμα ελέγχου, προφίλ και σπόρο (\eng{seed}). Οι δοκιμές συγκρίνουν τις μεταβάσεις των πινάκων με την άμεση εφαρμογή κινήσεων στο μοντέλο κυβιδίων. Ο Πίνακας~\ref{tab:generated-table-metadata} συνοψίζει τα μεταδεδομένα των παραγόμενων πινάκων.

\begin{table}[htbp]\centering
{\small
\begin{tabular}{L{0.27\textwidth}L{0.18\textwidth}rrr}
\hline
Table & Kind & Domain & Entries & Moves \\
\hline
CO & move\_table & 2187 & 39366 & 18 \\
CO & pruning\_table & 2187 & 2187 & 18 \\
EO & move\_table & 2048 & 36864 & 18 \\
EO & pruning\_table & 2048 & 2048 & 18 \\
UD slice & move\_table & 495 & 8910 & 18 \\
UD slice & pruning\_table & 495 & 495 & 18 \\
P2 corner perm. & move\_table & 40320 & 403200 & 10 \\
P2 corner perm. & pruning\_table & 40320 & 40320 & 10 \\
P2 U/D edge perm. & move\_table & 40320 & 403200 & 10 \\
P2 U/D edge perm. & pruning\_table & 40320 & 40320 & 10 \\
P2 slice edge perm. & move\_table & 24 & 240 & 10 \\
P2 slice edge perm. & pruning\_table & 24 & 24 & 10 \\
\hline
\end{tabular}
}

\caption{Μεταδεδομένα των παραγόμενων πινάκων κινήσεων και αποκοπής ανά συντεταγμένη: μέγεθος πεδίου τιμών, πλήθος εγγραφών και πλήθος επιτρεπτών κινήσεων, όπως καταγράφονται στο μανιφέστο (\eng{manifest}) \code{table_manifest_thesis_seed_2026.json}.}
\label{tab:generated-table-metadata}
\end{table}

Η παραγωγή των πινάκων είναι ντετερμινιστική. Για κάθε συντεταγμένη η γεννήτρια διατρέχει όλες τις τιμές του πεδίου ορισμού, αποκωδικοποιεί αντιπροσωπευτική κατάσταση, εφαρμόζει κάθε επιτρεπτή κίνηση και κωδικοποιεί ξανά το αποτέλεσμα. Ο ίδιος πίνακας κινήσεων χρησιμοποιείται έπειτα για BFS στον χώρο της προβολής, ώστε να παραχθεί ο αντίστοιχος πίνακας αποκοπής. Η διαδικασία γράφει τόσο τα ωμά αρχεία πινάκων στο \code{data/generated/} όσο και αντίγραφο μεταδεδομένων στο \code{results/processed/}. Η ύπαρξη δύο τοποθεσιών έχει πρακτικό σκοπό: τα αρχεία πινάκων είναι δεδομένα των επιλυτών, ενώ τα επεξεργασμένα μεταδεδομένα είναι δεδομένα τεκμηρίωσης των αποτελεσμάτων.

Τα αθροίσματα ελέγχου αποθηκεύονται επειδή οι πίνακες αλλάζουν αν αλλάξει ο κώδικας συντεταγμένων ή η σύμβαση κινήσεων. Το μανιφέστο κρατά επίσης το πεδίο \code{source_state}, δηλαδή σύντομη πληροφορία Git με την υποβολή (\eng{commit}) προέλευσης και ένδειξη καθαρότητας του δέντρου εργασίας. Τα τεκμήρια της εργασίας φέρουν σφραγίδα καθαρής, καταχωρισμένης κατάστασης πηγαίου κώδικα: το μανιφέστο πινάκων καταγράφει την υποβολή \code{dafe9ca} με ένδειξη \code{dirty=false}, ενώ τα μεταδεδομένα του γωνιακού πίνακα προτύπων καταγράφουν την υποβολή \code{f8d9a12}, επίσης με \code{dirty=false}. Έτσι κάθε τεκμήριο αντιστοιχίζεται σε συγκεκριμένη, αναπαραγώγιμη κατάσταση του πηγαίου κώδικα. Μοναδική εξαίρεση αποτελούν τα μεταδεδομένα του πίνακα \code{h48h7}, τα οποία διατηρούνται ως έχουν· η περίπτωση αυτή συζητείται στο Κεφάλαιο~\ref{chap:validation}.

Για τις μικρές προβολές χρησιμοποιείται JSON, επειδή η αναγνωσιμότητα και τα αθροίσματα ελέγχου είναι σημαντικότερα από τη μέγιστη συμπίεση. Για τον γωνιακό πίνακα προτύπων του 3×3×3, όμως, το JSON δεν θα ήταν κατάλληλο. Η εγγενής γεννήτρια στο \path{native/corner_pdb/} παράγει πλήρη πίνακα \(8!\cdot3^7\) καταστάσεων σε δυαδική μορφή, με επικεφαλίδα και ένα byte ανά απόσταση. Το πρόγραμμα \code{generate_corner_pdb.py} μεταγλωττίζει τη γεννήτρια C++, εκτελεί BFS πάνω στην προβολή των γωνιών, γράφει το δυαδικό τεκμήριο και αποθηκεύει μεταδεδομένα με άθροισμα ελέγχου. Η πλευρά Python δεν αντιγράφει τον πίνακα σε λίστα· τον ανοίγει με \code{mmap} μέσω του \code{tables/corner_pdb.py} και διαβάζει μόνο το byte της ζητούμενης προβεβλημένης κατάστασης. Τα μεταδεδομένα του πίνακα συνοψίζονται στον Πίνακα~\ref{tab:corner-pdb-metadata}.

\begin{table}[htbp]\centering
{\small
\begin{tabular}{lr}
\hline
Metric & Value \\
\hline
Projected states & 88179840 \\
Visited states & 88179840 \\
Complete & True \\
Maximum distance & 11 \\
Binary size bytes & 88179896 \\
Runtime seconds & 9.880 \\
Distribution buckets & 12 \\
\hline
\end{tabular}
}

\caption{Μεταδεδομένα του γωνιακού πίνακα προτύπων του 3×3×3: πλήθος προβεβλημένων καταστάσεων, πληρότητα, μέγιστη απόσταση, μέγεθος δυαδικού αρχείου και χρόνος παραγωγής, από το τεκμήριο \code{corner_pdb_metadata_seed_2026_thesis.json}.}
\label{tab:corner-pdb-metadata}
\end{table}

Ο γωνιακός πίνακας είναι πλήρης για την προβολή γωνιών του 3×3×3, όχι για ολόκληρο τον κύβο. Είναι ωστόσο πραγματικός πίνακας προτύπων του 3×3×3 και όχι υποκατάστατο από τον κύβο 2×2×2. Η τιμή του αποτελεί παραδεκτό κάτω φράγμα για το πλήρες πρόβλημα, επειδή κάθε λύση του πλήρους κύβου λύνει και την προβεβλημένη κατάσταση των γωνιών. Η IDA* χρησιμοποιεί το μέγιστο ανάμεσα στον γωνιακό πίνακα, στους μικρούς προβολικούς πίνακες και στο φράγμα λανθασμένων κυβιδίων.

Η επέκταση των πινάκων προτύπων στις ακμές γίνεται με οκτώ εξακμικές προβολές, οργανωμένες ως συμπληρωματικές και μικτές διαμερίσεις των ακμών. Κάθε προβολή έχει \(\binom{12}{6}\cdot6!\cdot2^6=42\,577\,920\) καταστάσεις και αποθηκεύεται σε δυαδικό τεκμήριο ενός byte ανά απόσταση. Η υλοποίηση συνδυάζει τους πίνακες με μέγιστο και όχι με άθροισμα, επειδή το άθροισμα επικαλυπτόμενων πινάκων πλήρους κόστους δεν είναι παραδεκτό χωρίς κατάτμηση κόστους (\eng{cost partitioning}) \cite{felner2004additivePDB}. Έτσι το φράγμα παραμένει ασφαλές για βέλτιστη αναζήτηση. Το συνοδευτικό πείραμα κάλυψης συγκρίνει το αρχικό σύνολο τεσσάρων προβολών με το σύνολο των οκτώ προβολών σε ντετερμινιστικό δείγμα. Καταγράφει μόνο ενίσχυση ή ισότητα του παραδεκτού φράγματος των ακμών· τα αποτελέσματα συνοψίζονται στον Πίνακα~\ref{tab:edge-pdb-coverage}. Επιπλέον, η εγγενής γεννήτρια υποστηρίζει πίνακες με κόστος τελεστών 0/1. Με αυτήν παρήχθησαν δύο πλήρεις προσθετικοί πίνακες ακμών (CPDB) με κατάτμηση κόστους, για τις ομάδες κινήσεων \(U/R/F\) και \(D/L/B\). Τα δύο αυτά τεκμήρια καλύπτουν 85\,155\,840 προβεβλημένες καταστάσεις. Φορτώνονται από τον εγγενή επιλυτή με τη χωριστή επιλογή \code{--additive-edge-pdb}, ώστε να μη συγχέονται με τους πίνακες πλήρους κόστους. Τα μεταδεδομένα των οκτώ εξακμικών πινάκων συνοψίζονται στον Πίνακα~\ref{tab:edge-pdb-metadata}.

\begin{table}[htbp]\centering
{\small
\begin{tabular}{lrrrrr}
\hline
Subset & States & Complete & Max dist. & Bytes & Seconds \\
\hline
0,1,2,3,4,5 & 42577920 & True & 10 & 42577992 & 283.784 \\
6,7,8,9,10,11 & 42577920 & True & 10 & 42577992 & 326.238 \\
0,2,4,6,8,10 & 42577920 & True & 10 & 42577992 & 317.767 \\
1,3,5,7,9,11 & 42577920 & True & 10 & 42577992 & 442.093 \\
0,1,4,5,8,9 & 42577920 & True & 10 & 42577992 & 305.119 \\
2,3,6,7,10,11 & 42577920 & True & 10 & 42577992 & 190.947 \\
0,3,5,6,8,11 & 42577920 & True & 10 & 42577992 & 354.691 \\
1,2,4,7,9,10 & 42577920 & True & 10 & 42577992 & 345.839 \\
\hline
\end{tabular}
}

\caption{Μεταδεδομένα των οκτώ εξακμικών πινάκων προτύπων του 3×3×3 ανά υποσύνολο ακμών: πλήθος καταστάσεων, πληρότητα, μέγιστη απόσταση, μέγεθος και χρόνος παραγωγής, από τα τεκμήρια \code{edge_pdb_metadata_subset_*.json}.}
\label{tab:edge-pdb-metadata}
\end{table}

\begin{table}[htbp]\centering
\input{tables/edge_pdb_coverage_expanded8.tex}
\caption{Σύγκριση κάλυψης του παραδεκτού φράγματος ακμών μεταξύ του αρχικού συνόλου τεσσάρων και του εκτεταμένου συνόλου οκτώ εξακμικών προβολών, σε ντετερμινιστικό δείγμα ανά βάθος ανακατέματος (\eng{scramble}), από το τεκμήριο \code{edge_pdb_coverage_seed_2026_thesis_expanded8.json}.}
\label{tab:edge-pdb-coverage}
\end{table}

Το πρόγραμμα \code{scripts/compare_heuristics.py} παράγει χωριστό τεκμήριο για τη σύνθεση της ευρετικής της IDA*. Το τεκμήριο συγκρίνει το φράγμα λανθασμένων κυβιδίων, τα μικρά προβολικά φράγματα αποκοπής, τον γωνιακό πίνακα προτύπων, τους πίνακες ακμών, το τελικό μέγιστο και, όταν υπάρχει, το χωριστό προσθετικό άθροισμα CPDB. Οι ρηχές γραμμές ελέγχονται με ακριβή απόσταση BFS, ενώ οι βαθύτερες γραμμές χρησιμοποιούνται μόνο για σύγκριση κάτω φραγμάτων. Τα αποτελέσματα της σύγκρισης παρουσιάζονται στο Κεφάλαιο~\ref{chap:experiments}.

Η εγγενής IDA* πλήρους κύβου στο \path{native/optimal_solver/} φορτώνει τον γωνιακό πίνακα προτύπων και όλους τους διαθέσιμους πίνακες ακμών. Η βασική ευρετική της είναι το μέγιστο των προβολικών αποστάσεων και επομένως δεν υπερεκτιμά. Όταν δοθεί συμβατό προσθετικό σύνολο με την επιλογή \code{--additive-edge-pdb}, ο επιλυτής αθροίζει τις αντίστοιχες αποστάσεις κόστους 0/1 και λαμβάνει το μέγιστο ανάμεσα στο άθροισμα αυτό και στα φράγματα πλήρους κόστους. Όταν η αναζήτηση ολοκληρώνεται και επιστρέφει \code{exact}, η λύση είναι βέλτιστη για τη συγκεκριμένη κατάσταση κυβιδίων στη μετρική HTM. Όταν σταματά από όριο χρόνου, βάθους ή κόμβων, το αποτέλεσμα παραμένει \code{timeout} ή κάτω φράγμα και δεν μετατρέπεται σε απόσταση.

Ειδική περίπτωση αποτελεί η αποκοπή κινήσεων ρίζας. Όταν είναι ενεργή η επιλογή \code{--root-move-mask}, ο εγγενής επιλυτής δεν επιστρέφει ανεπιφύλακτο \code{exact}, αλλά τη διακριτή ετικέτα \code{exact_under_root_mask}, επειδή η πληρότητα της απόδειξης εξαρτάται από την ορθότητα της μάσκας. Μόνο το περίβλημα Python αναβαθμίζει την ετικέτα σε \code{exact}, και μόνο όταν έχει παραγάγει το ίδιο τη μάσκα από τον πιστοποιημένο σταθεροποιητή συμμετρίας 48 στοιχείων της αρχικής κατάστασης· τότε η βελτιστότητα στο περιορισμένο δέντρο συνεπάγεται ανεπιφύλακτη βελτιστότητα. Η αναβάθμιση καταγράφεται ρητά στις σημειώσεις του αποτελέσματος, ενώ κάθε άλλο αποτέλεσμα \code{exact_under_root_mask} διατηρείται ως υπό όρους και δεν χαρακτηρίζεται επαληθευμένο.

\section{Εγγενές υποσύστημα H48}
\label{sec:implementation-h48}
Μετά την αρχική εγγενή υλοποίηση IDA* προστέθηκε δεύτερη, πιο εξειδικευμένη μηχανή ακριβούς αναζήτησης, βασισμένη στη δημόσια σχεδίαση H48 του \code{nissy-core}. Το υποσύστημα H48 (\eng{backend}) δεν είναι απλή κλήση σε εγκατεστημένο εκτελέσιμο. Η υλοποίηση περιλαμβάνει ενσωματωμένο στιγμιότυπο του \code{nissy-core} (άδεια GPL-3.0-or-later), περίβλημα C στο \path{native/h48_backend/h48_backend.c}, στρώση μεταγλώττισης και παραγωγής σε Python στο \code{tables/h48.py} και περίβλημα επιλυτή στο \code{solvers/h48_native.py}. Η βιβλιογραφική και τεχνική καταγωγή του υποσυστήματος ανήκει στη γραμμή εργασίας των Nissy/H48, nxopt και του βέλτιστου επιλυτή του Kociemba \cite{trontoNissyCoreH48,rokickiNxopt,kociembaOptimalSolver}.

Το όριο άδειας, προέλευσης και πατρότητας καταγράφεται ρητά και εκτός του κειμένου της εργασίας. Το αρχείο \path{THIRD_PARTY_NOTICES.md} παραπέμπει στην τοπική άδεια του ενσωματωμένου στιγμιότυπου,\footnote{\path{native/h48_backend/third_party/nissy_core/LICENSE}} στη διεύθυνση του ανάντη (\eng{upstream}) \code{nissy-core}, στην υποβολή που καταγράφουν ο κώδικας και τα παραγόμενα μεταδεδομένα, καθώς και στα δημόσια τεκμήρια πινάκων Nissy/RubikOptimal, τα οποία χρησιμοποιούνται μόνο ως εξωτερική ένδειξη. Η μηχανή H48/Nissy/RubikOptimal δεν αποτελεί πρωτότυπη αλγοριθμική συνεισφορά της εργασίας. Ο χαρακτηρισμός του υποσυστήματος ως ενσωματωμένου σημαίνει ότι η υλοποίηση μπορεί να το μεταγλωττίσει, να παράγει και να ελέγξει πίνακες και να καταγράψει αναπαραγώγιμα τεκμήρια· δεν σημαίνει ότι ο αλγόριθμος H48, ο πηγαίος κώδικας του \code{nissy-core} ή οι πίνακες H48 έχουν πατρότητα από τον συγγραφέα. Στους ισχυρισμούς συνεισφοράς, η μηχανή αυτή αναφέρεται πάντα χωριστά από την εγγενή ευρετική υποδομή Korf/IDA* και από τις εκπαιδευτικές υλοποιήσεις Kociemba και Thistlethwaite.

Η παραγωγή του πίνακα H48 είναι πλήρως αυτοματοποιημένη και αναπαραγώγιμη ως διαδικασία· ο διατηρούμενος πίνακας \code{h48h7}, ωστόσο, δεν αναπαράγεται πανομοιότυπος byte προς byte, για τους λόγους που τεκμηριώνονται στο Κεφάλαιο~\ref{chap:validation}. Το πρόγραμμα \code{scripts/generate_h48_tables.py} μεταγλωττίζει το περίβλημα, καλεί τη γεννήτρια H48, αποθηκεύει τον δυαδικό πίνακα κάτω από το \path{data/generated/h48/}, υπολογίζει SHA-256 και γράφει μεταδεδομένα στο \path{results/processed/}. Το προφίλ της εργασίας περιέχει δύο επίπεδα: τον μικρό πίνακα \code{h48h0} για γρήγορη αναπαραγωγή και τον πίνακα \code{h48h7} επιπέδου μαντείου (\eng{oracle}), ο οποίος αντιστοιχεί στο ψευδώνυμο \code{optimal} του \code{nissy-core}. Το τεκμήριο \code{h48h7} έχει μέγεθος 3\,793\,842\,344 bytes. Παράχθηκε μετά από διόρθωση της τοπικής ανίχνευσης \code{pthread} στο ενσωματωμένο περίβλημα, ώστε η γεννήτρια να χρησιμοποιεί πραγματικά νήματα εργασίας σε macOS. Ο Πίνακας~\ref{tab:h48-metadata} συνοψίζει τα παραγόμενα τεκμήρια H48.

\begin{table}[htbp]\centering
{\small
\begin{tabular}{llrll}
\hline
Solver & Profile & Bytes & SHA-256 prefix & Status \\
\hline
h48h0 & quick & 31683944 & 46dcee74f59c & generated \\
h48h0 & thesis & 31683944 & 46dcee74f59c & generated \\
h48h7 & thesis & 3793842344 & eacc59ecac61 & generated \\
\hline
\end{tabular}
}

\caption{Μεταδεδομένα των παραγόμενων πινάκων αποκοπής H48 ανά επίπεδο και προφίλ: μέγεθος σε bytes, πρόθεμα SHA-256 και κατάσταση παραγωγής, από τα τεκμήρια \code{h48_metadata_seed_2026_*.json}.}
\label{tab:h48-metadata}
\end{table}

Η μηχανή H48 τηρεί αυστηρό συμβόλαιο αναφοράς. Επιστρέφει \code{exact} μόνο όταν το εγγενές υποσύστημα ολοκληρώσει την αναζήτηση και η πλευρά Python επαληθεύσει τη λύση στον τοπικό μηχανισμό επαλήθευσης κυβιδίων. Αν η θυγατρική διεργασία ξεπεράσει το καθορισμένο όριο χρόνου, η ετικέτα γίνεται \code{timeout}. Αν λείπει πίνακας ή αν η είσοδος αποτύχει στον έλεγχο φυσικής εγκυρότητας, η ετικέτα γίνεται \code{failed} ή η περίπτωση απορρίπτεται πριν από την κλήση του υποσυστήματος. Η υλοποίηση μετατρέπει απευθείας την τοπική κατάσταση (σε μορφή κυβιδίων ή ψηφίδων) στη συμπαγή συμβολοσειρά του \code{nissy-core}, με αντιστοίχιση των διαφορετικών διατάξεων γωνιών και ακμών. Το τεκμήριο καταπόνησης της μηχανής H48 δεν χρειάζεται να περάσει στο υποσύστημα την ακολουθία που παρήγαγε την κατάσταση: οι γραμμές των μη λυμένων καταστάσεων καταγράφουν \code{input_mode=cube_state} και \code{source_sequence_provided=False}.

Το τεκμήριο του \code{h48h7} επεκτείνει αυτή την απόδειξη σε 15 γραμμές άμεσης εισόδου κατάστασης (\eng{direct-state}), συμπεριλαμβανομένων δέκα ντετερμινιστικών περιπτώσεων βάθους 25. Το χωριστό πρόγραμμα πιστοποίησης προσθέτει τη λυμένη κατάσταση, μία ρηχή κατάσταση, μία ντετερμινιστική κατάσταση βάθους 25 και την τυπική κατάσταση \eng{superflip}, στην οποία όλες οι ακμές είναι ανεστραμμένες, με αναμενόμενη απόσταση 20 στη μετρική HTM. Η κλήση γίνεται στη συνάρτηση \code{nissy_solve} με \code{max_depth=20}, και το \code{nissy-core} τεκμηριώνει το H48 ως βέλτιστο επιλυτή HTM χωρίς προαπαιτούμενα. Αυτό ενισχύει τη σύμβαση αναφοράς αποκλειστικά ακριβών αποτελεσμάτων και την τεκμηρίωση των αποθηκευμένων γραμμών. Η εργασία ωστόσο εξακολουθεί να διαχωρίζει αυτή την αλγοριθμική σύμβαση και τις αποθηκευμένες ακριβείς γραμμές από ένα τυπικό θεώρημα χειρότερης περίπτωσης χρόνου εκτέλεσης για κάθε πιθανή κατάσταση 3×3×3.

Για επαναλαμβανόμενη χρήση του μαντείου προστέθηκε λειτουργία \code{solve-batch} στο εγγενές περίβλημα. Στη λειτουργία μεμονωμένης κλήσης κάθε κλήση ξεκινά νέα διεργασία, διαβάζει τον δυαδικό πίνακα \code{h48h7} και εκτελεί τον πλήρη έλεγχο \code{nissy_checkdata}. Στη λειτουργία δέσμης ο ίδιος πίνακας φορτώνεται και ελέγχεται μία φορά, και στη συνέχεια πολλές συμπαγείς συμβολοσειρές καταστάσεων λύνονται από την ίδια διεργασία. Το περίβλημα Python \code{solve_h48_native_batch} κρατά το αποτέλεσμα \code{exact} ή \code{timeout} ανά κατάσταση και επαληθεύει κάθε λύση όπως και στη μεμονωμένη κλήση. Η δημόσια εντολή \code{rubik-optimal oracle} χρησιμοποιεί αυτή τη λειτουργία δέσμης για είσοδο ανά γραμμή· παραδείγματα χρήσης της δίνονται στο Κεφάλαιο~\ref{chap:usage}.

Επιπλέον, το πακέτο εκθέτει τη διεπαφή προγραμματισμού \code{FastOptimalOracle} στο \path{src/rubik_optimal/oracle.py}. Η κλάση αυτή είναι η ολοκληρωμένη μορφή του βελτιστοποιημένου μαντείου: επιλέγει ως προεπιλογή το \code{h48h7}, κρατά μόνιμη εγγενή διεργασία και χρησιμοποιεί τα έμπιστα μεταδεδομένα των παραγόμενων πινάκων, ώστε να αποφεύγει την επαναλαμβανόμενη πλήρη σάρωση του πίνακα. Επιπρόσθετα, απορρίπτει μη φυσικές καταστάσεις 3×3×3 πριν από την κλήση του υποσυστήματος και επαληθεύει κάθε επιστρεφόμενη λύση στον ανεξάρτητο μηχανισμό επαλήθευσης. Η λυμένη κατάσταση επιστρέφεται αμέσως χωρίς εγγενή διεργασία. Οι μεταβλητές περιβάλλοντος \code{RUBIK_OPTIMAL_H48_THREADS} και \code{RUBIK_OPTIMAL_THREADS} επιτρέπουν πιο συντηρητική χρήση CPU σε φορτωμένο μηχάνημα, χωρίς αλλαγή της υλοποίησης του επιλυτή.

Η σημασιολογία των χρόνων διαφέρει ανά λειτουργία. Για αποτελέσματα μεμονωμένης κλήσης το πεδίο \code{runtime_seconds} μετρά το πλήρες χρονικό διάστημα της θυγατρικής διεργασίας Python, ενώ το \code{backend_solve_seconds} παραμένει στις σημειώσεις ως ο καθαρός χρόνος αναζήτησης του \code{nissy-core}. Για αποτελέσματα δέσμης το \code{runtime_seconds} είναι ο χρόνος επίλυσης του υποσυστήματος ανά περίπτωση, ενώ ο κοινός χρόνος \code{batch_wall_seconds} καταγράφεται στις σημειώσεις, ώστε να μη συγχέεται η ρυθμαπόδοση με τη δυσκολία χειρότερης περίπτωσης μιας μεμονωμένης κατάστασης.

Η εγγενής φόρτωση του πίνακα H48 χρησιμοποιεί \code{mmap} όταν η στοίχιση μνήμης το επιτρέπει και καταφεύγει σε στοιχισμένη μνήμη σωρού μόνο αν χρειαστεί. Για παραγόμενα τεκμήρια με επαληθευμένο άθροισμα ελέγχου υπάρχει η προαιρετική επιλογή \code{--h48-trusted-table}, η οποία παραλείπει την πλήρη σάρωση \code{nissy_checkdata} σε κάθε κλήση. Η επιλογή αυτή δεν είναι νέος αλγόριθμος αναζήτησης· αφαιρεί επαναλαμβανόμενο κόστος επικύρωσης, όταν ο ίδιος πίνακας έχει ήδη παραχθεί και ταυτοποιηθεί από μεταδεδομένα και άθροισμα ελέγχου. Για δύσκολες μεμονωμένες καταστάσεις υπάρχει επίσης η επιλογή \code{--h48-preload-table}, η οποία αγγίζει σειριακά τις σελίδες του απεικονισμένου πίνακα πριν από την αναζήτηση, ώστε να μειώνονται τα ψυχρά σφάλματα σελίδας. Οι σημειώσεις των αποτελεσμάτων καταγράφουν τα πεδία \code{table_check}, \code{table_storage} και \code{table_preload}, ώστε η ανάλυση του Κεφαλαίου~\ref{chap:experiments} να ξεχωρίζει τον γρήγορο έλεγχο πίνακα από πραγματική επιτάχυνση της ίδιας της αναζήτησης H48.

\section{Υλοποίηση IDA*}
\label{sec:implementation-ida}
Η IDA* βρίσκεται στο \code{search/ida_star.py}. Η συνάρτηση δέχεται αρχική κατάσταση, μέγιστο βάθος, όριο χρόνου, όριο κόμβων, ευρετική, προαιρετικό κατηγόρημα στόχου και προαιρετικό σύνολο κινήσεων. Μπορεί επίσης να ταξινομεί τα παιδιά με βάση την ευρετική τιμή, χρησιμοποιώντας τοπική κρυφή μνήμη της ευρετικής, ώστε να μην υπολογίζεται η ίδια προβολή επανειλημμένα. Η παραμετροποίηση αυτή επιτρέπει να χρησιμοποιείται η ίδια μηχανή τόσο για πλήρη επίλυση προς τη λυμένη κατάσταση όσο και για αναζητήσεις με περιορισμένο σύνολο κινήσεων. Η τετραφασική υλοποίηση Thistlethwaite δεν χρησιμοποιεί αυτή τη μηχανή, επειδή καθοδηγείται από ακριβείς φασικούς πίνακες αποστάσεων, όπως περιγράφεται στην Ενότητα~\ref{sec:implementation-thistlethwaite}.

Το αποτέλεσμα της IDA* περιλαμβάνει λύση ή \code{None}, πλήθος επεκταθέντων και παραχθέντων κόμβων, ετικέτα κατάστασης, χρόνο εκτέλεσης, αρχικό κάτω φράγμα και επεξηγηματικές σημειώσεις. Οι σημειώσεις βοηθούν ιδιαίτερα στις υπερβάσεις ορίων: μια διακοπή επειδή το κάτω φράγμα ξεπερνά το καθορισμένο βάθος είναι διαφορετική από διακοπή λόγω ορίου χρόνου ή κόμβων. Και στις δύο περιπτώσεις η ετικέτα είναι \code{timeout}, αλλά η σημείωση καθοδηγεί την ερμηνεία.

Η αναδρομική αναζήτηση χρησιμοποιεί \(f=g+h\). Αν το \(f\) ξεπερνά το τρέχον όριο, ο κλάδος κόβεται και επιστρέφει την τιμή που θα μπορούσε να γίνει το επόμενο όριο. Αν η κατάσταση ικανοποιεί το κατηγόρημα στόχου, επιστρέφεται το μονοπάτι. Αν συμπληρωθεί το μέγιστο βάθος, το όριο χρόνου ή το όριο κόμβων, ο κλάδος σταματά. Η υλοποίηση αποφεύγει διαδοχικές κινήσεις στην ίδια έδρα μέσω της βοηθητικής συνάρτησης \code{same_face}.

Η μηχανή IDA* δεν αποθηκεύει πίνακα ήδη εξερευνημένων καταστάσεων (\eng{transposition table}). Η επιλογή αυτή κρατά την υλοποίηση απλή και ελαφριά σε μνήμη, αλλά επηρεάζει την απόδοση: σε βαθύτερες αναζητήσεις οι ίδιες καταστάσεις μπορούν να εμφανιστούν ξανά από διαφορετικά μονοπάτια. Η προσθήκη τέτοιας αποκοπής θα ήταν πιθανή μελλοντική βελτίωση, αλλά θα έπρεπε να γίνει προσεκτικά, ώστε να μην αλλοιωθεί η σχέση με το επαναληπτικά αυξανόμενο όριο.

\section{Υλοποίηση Kociemba}
\label{sec:implementation-kociemba}
Η υλοποίηση Kociemba βρίσκεται στο \code{solvers/kociemba.py}. Η συνάρτηση \code{is_kociemba_phase1_goal} ορίζει την ενδιάμεση ομάδα με βάση τις συντεταγμένες CO, EO και UD-slice. Η \code{solve_kociemba_phase1} δεν επεκτείνει πλήρεις καταστάσεις του κύβου σε αυτό το στάδιο. Εκτελεί IDA* απευθείας στο γινόμενο των τριών παραγόμενων πινάκων κινήσεων CO, EO και UD-slice, με στόχο την τριάδα των λυμένων συντεταγμένων. Ως κάτω φράγμα χρησιμοποιείται από προεπιλογή ο ακριβής, συμμετρικά ανηγμένος πίνακας αποστάσεων φάσης~1 βάθους 12,\footnote{Δυαδικό τεκμήριο \path{data/generated/kociemba_phase1_sym_depth12.bin}.} ο οποίος κυριαρχεί σε καθεμία από τις τρεις μονοδιάστατες προβολές του ίδιου χώρου. Αν ο πίνακας αυτός δεν είναι διαθέσιμος ή δεν περνά τους ελέγχους ακεραιότητας, η υλοποίηση επιστρέφει στο μέγιστο των αντίστοιχων αποστάσεων αποκοπής. Οι σημειώσεις του αποτελέσματος καταγράφουν την επιλεγμένη ευρετική (\code{phase1_heuristic=sym_reduced_exact_depth12} ή \code{projection_max}).

Η υλοποίηση δεν δεσμεύεται στην πρώτη κατάσταση που ικανοποιεί τη φάση 1. Η \code{collect_kociemba_phase1_candidates} συλλέγει φραγμένο σύνολο υποψήφιων μεταβάσεων προς την ενδιάμεση ομάδα και ο επιλυτής τις ταξινομεί με βάση το κάτω φράγμα της φάσης 2. Έπειτα εφαρμόζει κάθε υποψήφια ακολουθία στην αρχική κατάσταση και εκτελεί τη \code{solve_kociemba_phase2} στη νέα κατάσταση. Η φάση 2 εκτελεί αντίστοιχη IDA* στον χώρο συντεταγμένων, πάνω στις φασικές προβολές μετάθεσης γωνιών, μετάθεσης ακμών U/D και μετάθεσης ακμών φέτας, με το περιορισμένο σύνολο κινήσεων \code{PHASE2_MOVES}, το οποίο αντιστοιχεί στα σύμβολα \(U,U',U2,D,D',D2,R2,L2,F2,B2\). Οι δύο φάσεις κρατούν ανά όριο μικρό πίνακα ήδη εξερευνημένων καταστάσεων, με κλειδί τη συντεταγμένη και την προηγούμενη έδρα, ώστε να μη διερευνούν ξανά την ίδια προβεβλημένη κατάσταση σε ίσο ή μεγαλύτερο βάθος. Οι πίνακες προθερμαίνονται πριν ξεκινήσει το μετρούμενο παράθυρο αναζήτησης, ώστε ο χρόνος της οκνηρής (\eng{lazy}) κατασκευής τους να μη συγχέεται με το όριο χρόνου της ίδιας της αναζήτησης.

Η τελική λύση είναι συνένωση της φάσης 1 και της φάσης 2 και επαληθεύεται ξανά στον πλήρη κύβο. Ανεξάρτητα από το αν οι φάσεις βρήκαν λύση, η ετικέτα δεν γίνεται \code{exact}: η μέθοδος είναι σταδιακή και περιορισμένης εμβέλειας, όχι πλήρης απόδειξη βελτιστότητας. Η υλοποίηση επιστρέφει \code{non_exact} για επαληθευμένη λύση, \code{lower_bound} όταν το παραδεκτό φράγμα ή το καθορισμένο βάθος αποκλείει ολοκλήρωση εντός των ορίων, και \code{timeout} όταν το όριο χρόνου ή κόμβων σταματήσει την αναζήτηση. Τα προεπιλεγμένα όρια βάθους ισούνται με τις φασικές διαμέτρους (12 για τη φάση 1 και 18 για τη φάση 2), ώστε καμία έγκυρη κατάσταση να μην αποκλείεται από τη διαμόρφωση. Μόνο τα όρια χρόνου και κόμβων φράσσουν την αναζήτηση. Ο καταμερισμός του χρονικού προϋπολογισμού εγγυάται ότι η φάση 2 λαμβάνει πάντα φραγμένη απόπειρα όταν υπάρχουν υποψήφιες μεταβάσεις, αντί να εξαντλείται όλος ο διαθέσιμος χρόνος στη συλλογή της φάσης 1. Με τις προεπιλογές αυτές η υλοποίηση λύνει τυπικά τυχαία ανακατέματα 18--20 κινήσεων, αλλά τα αποτελέσματα παραμένουν επαληθευμένες πρακτικές λύσεις και όχι αποδείξεις βέλτιστου μήκους.

\section{Συμμετρικά ανηγμένη φάση 1 και εγγενής απόδειξη βελτιστότητας}
\label{sec:implementation-sym-proof}
Πάνω στην υλοποίηση Kociemba χτίζεται μια εγγενής (χωρίς H48/Nissy) μηχανή απόδειξης βελτιστότητας, υλοποιημένη στο \path{native/kociemba_phase2_probe/}. Το κρίσιμο εργαλείο είναι ένας συμπιεσμένος, συμμετρικά ανηγμένος πίνακας αποκοπής της φάσης~1, στο πρότυπο των Kociemba/Cube~Explorer \cite{kociembaOptimalSolver,kociembaImplementationDetails}. Η συνδυασμένη συντεταγμένη \eng{FlipUDSlice} (προσανατολισμός ακμών \(\times\) συνδυασμός UD-slice, \(2048\cdot495=1\,013\,760\) τιμές) ανάγεται κατά τις 16 συμμετρίες ολόκληρου του κύβου που σταθεροποιούν τον άξονα U/D, δίνοντας τις 64\,430 κλάσεις ισοδυναμίας του Kociemba. Σε αντίθεση με μια αφελή παραγοντοποίηση, η \eng{FlipUDSlice} δεν αναλύεται σε ανεξάρτητες δράσεις προσανατολισμού και φέτας: η συζυγία συμμετρίας αναμειγνύει τον προσανατολισμό ακμών με τη μετάθεση ακμών, άρα η αναγωγή γίνεται στη συνδυασμένη συντεταγμένη. Όλη η αλγεβρική εργασία, ορθή και για τις ανακλάσεις, γίνεται στο επικυρωμένο επίπεδο συμμετρίας \code{src/rubik_optimal/symmetry.py} μέσω του \code{scripts/generate_phase1_sym_tables.py}, το οποίο μετασχηματίζει μόνο τους περίπου 64\,430 αντιπροσώπους. Για τις 2033 κλάσεις με μη τετριμμένο σταθεροποιητή, ο ανηγμένος προσανατολισμός γωνιών πρέπει να είναι το ελάχιστο πάνω στο σύμπλοκο (\eng{coset}) του σταθεροποιητή. Μια αυθαίρετη επιλογή έδινε μη παραδεκτή υπερεκτίμηση, η οποία εντοπίστηκε και διορθώθηκε. Ο εγγενής φορτωτής χτίζει με BFS τον πίνακα \(64\,430\times2187=140\,908\,410\) εγγραφών, πλήρη μέχρι τη διάμετρο~12 της φάσης~1. Ο πίνακας αυτός καταλαμβάνει περίπου 141~MB έναντι 2.07~GiB του ωμού πίνακα, μέγεθος κρίσιμο για το μηχάνημα των 16~GiB.

Η ορθότητα του συμμετρικού πίνακα κατοχυρώνεται με σύγκριση byte προς byte έναντι του έμπιστου ωμού πίνακα BFS της φάσης~1, σε όλες τις 959\,761\,462 γνωστές εγγραφές με απόσταση φάσης~1 έως 9 (όσες γεμίζει ο ωμός πίνακας βάθους~9), με μηδέν αναντιστοιχίες, όπως καταγράφεται στο αποθηκευμένο τεκμήριο επαλήθευσης.\footnote{Τεκμήριο \path{results/processed/native_sym_phase1_verify_seed_2026_thesis.json}.} Οι εγγραφές απόστασης 10--12 κληρονομούν την ορθότητα από τον ίδιο επικυρωμένο γεννήτορα BFS, αφού ο συμμετρικός πίνακας χτίζεται πλήρης μέχρι τη διάμετρο~12.

Πάνω σε αυτόν τον πίνακα προστίθεται το παραδεκτό φράγμα τριών αξόνων του Mike Reid \cite{reid1997optimal}: η απόσταση φάσης~1 υπολογίζεται για τον άξονα U/D και, μέσω συζυγίας με δύο σταθερές περιστροφές, για τους άξονες R/L και F/B, και λαμβάνεται το μέγιστο \(\max(p_{ud},p_{rl},p_{fb})\). Επειδή κάθε όρος είναι κάτω φράγμα της απόστασης προς τη λυμένη κατάσταση, το μέγιστό τους είναι επίσης παραδεκτό φράγμα. Η ορθότητα της συζυγίας επιβεβαιώνεται μέσω της ομομορφικής ιδιότητας \(\varphi(s\cdot m)\varphi^{-1} = (\varphi s\varphi^{-1})(\varphi m\varphi^{-1})\), ώστε οι συντεταγμένες των τριών αξόνων να ανανεώνονται σταδιακά. Επειδή το \eng{superflip} είναι αναλλοίωτο σε όλες τις 48 συμμετρίες, η πρώτη κίνηση περιορίζεται στους αντιπροσώπους τροχιάς \code{U} και \code{U2} (αποκοπή ρίζας από τον σταθεροποιητή), μια ασφαλή ως προς τη βελτιστότητα μείωση.

Η ίδια η απόδειξη εκτελείται από έναν IDA* ενιαίου ορίου τύπου Reid (\code{--mode optimal-ida}), ο οποίος αναζητά απευθείας τη λυμένη κατάσταση με το πλήρες σύνολο κινήσεων, χρησιμοποιώντας το παραπάνω φράγμα τριών αξόνων ως ευρετική. Για να αποδειχθεί η απουσία λύσης μήκους \(\le T\), εκτελείται μία αναζήτηση με όριο \(T\): αν εξαντληθεί χωρίς να φτάσει σε λυμένο φύλλο, δεν υπάρχει λύση τέτοιου μήκους. Η ορθότητα της διαδικασίας ελέγχεται σε ρηχές περιπτώσεις, όπου η ίδια αναζήτηση βρίσκει βέλτιστη λύση στη σωστή απόσταση και αποδεικνύει την απουσία λύσης ένα βήμα χαμηλότερα. Το πλήθος των κόμβων είναι ντετερμινιστικό, ανεξάρτητο από νήματα ή κατάτμηση εργασιών. Τα αριθμητικά αποτελέσματα της απόδειξης για το \eng{superflip} παρουσιάζονται στο Κεφάλαιο~\ref{chap:experiments}.

\section{Υλοποίηση Thistlethwaite}
\label{sec:implementation-thistlethwaite}
Η εγγενής υλοποίηση Thistlethwaite βρίσκεται στο \path{solvers/thistlethwaite.py} και στηρίζεται στο \path{tables/thistlethwaite_tables.py}, το οποίο κατασκευάζει, αποθηκεύει και επαναφορτώνει τέσσερις ακριβείς πίνακες αποστάσεων BFS, έναν για κάθε μείωση υποομάδας. Ο πίνακας \(G_0 \rightarrow G_1\) έχει 2048 εγγραφές πάνω στη συντεταγμένη προσανατολισμού ακμών, με το πλήρες σύνολο των 18 κινήσεων HTM και μέγιστη απόσταση 7. Ο πίνακας \(G_1 \rightarrow G_2\) έχει \(2187\cdot495=1\,082\,565\) εγγραφές πάνω στο γινόμενο προσανατολισμού γωνιών και συνδυασμού UD-slice, με το σύνολο κινήσεων \code{G1_MOVES} και μέγιστη απόσταση 10· ο στόχος του συμπίπτει με την κατάσταση περιορισμών της πρώτης φάσης Kociemba. Ο πίνακας \(G_2 \rightarrow G_3\) ορίζεται στο γινόμενο των 420 αριστερών συμπλόκων μετάθεσης γωνιών και των 140 συμπλόκων μετάθεσης ακμών ως προς τη δράση της τετραγωνικής (\eng{square}) υποομάδας: από τις 58\,800 εγγραφές, οι 29\,400 αντιστοιχούν σε προσβάσιμα σύμπλοκα, με μέγιστη απόσταση 13. Ο πίνακας \(G_3 \rightarrow G_4\) καλύπτει ολόκληρη την τετραγωνική υποομάδα με \(96\cdot6912=663\,552\) εγγραφές και μέγιστη απόσταση 15. Οι πίνακες γράφονται στον δίσκο ως δυαδικά αρχεία ενός byte ανά εγγραφή, συνολικού μεγέθους περίπου 1.8~MB. Συνοδεύονται από αθροίσματα ελέγχου SHA-256 και μανιφέστο στο \path{data/generated/thistlethwaite/}, ώστε κατεστραμμένη ή ασύμβατη κρυφή μνήμη να εντοπίζεται και να αναπαράγεται, αντί να οδηγεί σε σιωπηλά λανθασμένη επίλυση. Οι απεικονίσεις κατάστασης σε δείκτη συμπλόκου ανακατασκευάζονται ντετερμινιστικά κατά τη φόρτωση και δένονται με τους πίνακες μέσω των αθροισμάτων ελέγχου του μανιφέστου.

Η επίλυση κάθε φάσης είναι άπληστη κάθοδος καθοδηγούμενη από τον αντίστοιχο πίνακα (\code{_descend_phase}): σε κάθε βήμα δοκιμάζονται οι επιτρεπτές κινήσεις της φάσης και επιλέγεται κίνηση που μειώνει αυστηρά την αποθηκευμένη απόσταση. Επειδή ο πίνακας περιέχει την πραγματική απόσταση, τέτοια κίνηση υπάρχει πάντα μέχρι την απόσταση μηδέν, οπότε κάθε φάση τερματίζει εγγυημένα σε ακριβώς τόσες κινήσεις όση η αρχική εγγραφή του πίνακα. Σε αυτή την υλοποίηση δεν υπάρχει IDA*, συλλογή υποψήφιων μεταβάσεων, όριο χρόνου ή όριο κόμβων· οι παλαιότερες παράμετροι βάθους και ορίων διατηρούνται στη δημόσια υπογραφή μόνο για συμβατότητα με υπάρχουσες ροές μετρήσεων και αγνοούνται. Οι σημειώσεις του αποτελέσματος καταγράφουν τα τέσσερα φασικά μήκη και το συνολικό μέγεθος των αποθηκευμένων πινάκων. Τα κατηγορήματα ένταξης \code{is_g2_subgroup} και \code{is_g3_square_group} διατηρούνται ως ανεξάρτητοι έλεγχοι: το δεύτερο παράγει τα σύνολα των 96 μεταθέσεων γωνιών και των 6912 μεταθέσεων ακμών από τους ίδιους τους γεννήτορες μισών στροφών, ώστε να καλύπτεται και η πληροφορία συστροφής τετράδων (\eng{tetrad twist}) της τετραγωνικής υποομάδας.

Η υλοποίηση δεν χρησιμοποιεί τους ιστορικούς στατικούς πίνακες χειρισμών (\eng{maneuver tables})· τους αντικαθιστά με ακριβείς πίνακες αποστάσεων ανά φάση, και τα μέγιστα φασικά βάθη 7/10/13/15 που προκύπτουν συμπίπτουν με τα δημοσιευμένα φασικά φράγματα της μεθόδου \cite{scherphuisThistlethwaite}. Η τετραφασική αλυσίδα δεν συγχέεται με βέλτιστη αναζήτηση: κάθε λύση που περνά την ανεξάρτητη επαλήθευση χαρακτηρίζεται \code{non_exact}, όχι \code{exact}, ενώ αποτυχία της επαλήθευσης θα κατέγραφε \code{failed}. Η \code{solve_thistlethwaite_native_scoped} παραμένει ως συμβατό περίβλημα για τις υπάρχουσες ροές μετρήσεων και καλεί την ίδια τετραφασική υλοποίηση με το παλαιό όνομα επιλυτή.

\section{Κύβος 2×2×2}
\label{sec:implementation-pocket}
Ο φάκελος \code{pocket/} περιέχει το κανονικοποιημένο μοντέλο του κύβου 2×2×2 (\eng{Pocket Cube}). Η γωνία DBL παραμένει σταθερή, άρα ο χώρος καταστάσεων έχει επτά μετατιθέμενες γωνίες και έξι ανεξάρτητους προσανατολισμούς. Η πλήρης BFS δεν αποθηκεύει γονείς για όλες τις καταστάσεις· αποθηκεύει μόνο αποστάσεις, ώστε να παραχθεί η κατανομή. Για μεμονωμένες αντιπροσωπευτικές περιπτώσεις χρησιμοποιείται BFS με γονείς, ώστε να ανακτηθεί συγκεκριμένη βέλτιστη λύση.

Η μονάδα του 2×2×2 ορίζει την κατάσταση, τις κινήσεις και τους πίνακες μετάβασης. Το αρχείο \code{pocket/optimal.py} περιέχει δύο λειτουργίες: πλήρη κατανομή αποστάσεων και μεμονωμένη ανάκτηση βέλτιστης λύσης. Η πρώτη χρειάζεται να αποθηκεύσει μόνο απόσταση ανά συντεταγμένη. Η δεύτερη χρειάζεται πληροφορία μονοπατιού για τη συγκεκριμένη αρχική κατάσταση που εξετάζεται.

Η πλήρης κατανομή χρησιμοποιεί πίνακα μικρών ακεραίων αρχικοποιημένο με \(-1\). Η λυμένη κατάσταση έχει απόσταση 0. Η ουρά της BFS περιέχει συντεταγμένες. Κάθε άγνωστη γειτονική συντεταγμένη παίρνει απόσταση \(d+1\), αυξάνει την κατανομή και εισέρχεται στην ουρά. Όταν η ουρά εξαντληθεί, η κατανομή είναι πλήρης αν το πλήθος των επισκέψεων ισούται με τον αναμενόμενο χώρο των \(3\,674\,160\) καταστάσεων. Ο μηχανισμός επαλήθευσης ελέγχει ακριβώς αυτή τη συνθήκη στην αποθηκευμένη σύνοψη αποτελεσμάτων.

\section{Περιβάλλον γραμμής εντολών και σενάρια παραγωγής}
\label{sec:implementation-cli}
Το περιβάλλον γραμμής εντολών συγκεντρώνει τις βασικές εντολές για επίλυση, αναγνώριση απόστασης, χρήση του μαντείου σε δέσμες, μετρήσεις και παραγωγή πινάκων. Τα σενάρια (\eng{scripts}) του φακέλου \code{scripts/} παρέχουν σταθερές ροές παραγωγής των τεκμηρίων της εργασίας. Η διάκριση είναι πρακτική: το περιβάλλον γραμμής εντολών απευθύνεται σε χρήση από άνθρωπο ή σε γρήγορους ελέγχους, ενώ τα σενάρια παράγουν αναπαραγώγιμα τεκμήρια με συγκεκριμένα ονόματα αρχείων. Αναλυτικά παραδείγματα χρήσης των εντολών δίνονται στο Κεφάλαιο~\ref{chap:usage}, ενώ η πλήρης αναφορά των εντολών βρίσκεται στο Παράρτημα~\ref{app:cli}.

Οι βασικές ροές παραγωγής είναι οι εξής:

\begin{itemize}
\item \code{generate_tables.py}: πίνακες συντεταγμένων και μεταδεδομένα,
\item \code{generate_corner_pdb.py}: εγγενής παραγωγή του δυαδικού γωνιακού πίνακα προτύπων του 3×3×3,
\item \code{generate_edge_pdb.py}: εγγενής παραγωγή των δυαδικών εξακμικών πινάκων προτύπων του 3×3×3,
\item \code{generate_h48_tables.py}: εγγενής παραγωγή πινάκων αποκοπής H48 μέσω του ενσωματωμένου \code{nissy-core},
\item \code{run_h48_oracle_certification.py}: σώμα πιστοποίησης H48 άμεσης εισόδου κατάστασης, με το \eng{superflip},
\item \code{benchmark_h48_batch_overhead.py}: μέτρηση του αποσβεσμένου κόστους φόρτωσης του πίνακα H48 σε λειτουργία δέσμης,
\item \code{benchmark_h48_trusted_table.py}: μέτρηση του κόστους επικύρωσης του πίνακα H48 με και χωρίς έμπιστο πίνακα,
\item \code{run_h48_oracle_cli.py}: παραγωγή τεκμηρίων για τη δημόσια εντολή \code{rubik-optimal oracle},
\item \code{run_benchmarks.py}: ωμά αρχεία JSONL, επεξεργασμένα CSV και συνόψεις JSON,
\item \code{run_3x3_optimal.py}: τεκμήρια εγγενούς βέλτιστης αναζήτησης με πίνακες προτύπων γωνιών και ακμών ή με το υποσύστημα H48,
\item \code{generate_figures.py}: πίνακες LaTeX και αρχεία δεδομένων σχημάτων από τις μετρήσεις,
\item \code{generate_pocket_cube.py}: πλήρης κατανομή αποστάσεων του 2×2×2,
\item \code{verify_results.py} και \code{thesis_audit.py}: αυτοματοποιημένοι έλεγχοι ακεραιότητας των αποτελεσμάτων και των τεκμηρίων· ο ρόλος τους περιγράφεται στο Κεφάλαιο~\ref{chap:design} και το σχήμα που ελέγχουν τεκμηριώνεται στο Παράρτημα~\ref{app:schema}.
\end{itemize}

Η ύπαρξη πολλών σεναρίων δεν είναι τυχαία: κάθε σενάριο αντιστοιχεί σε διαφορετικό στάδιο παραγωγής, ώστε μια αλλαγή σε ένα επίπεδο να απαιτεί επανεκτέλεση μόνο των σταδίων που επηρεάζει. Ο διαχωρισμός αυτός αποτρέπει τόσο τις περιττές ακριβές επανεκτελέσεις όσο και την παρουσίαση παρωχημένου τεκμηρίου ως τρέχοντος.

\section{Δοκιμές}
\label{sec:implementation-tests}
Η δοκιμαστική σουίτα καλύπτει τον αναλυτή κινήσεων, τις κινήσεις κυβιδίων, τη μετατροπή μεταξύ μορφής κυβιδίων και μορφής ψηφίδων, τη φυσική εγκυρότητα, τις βοηθητικές συναρτήσεις αναζήτησης, τη συμπεριφορά των επιλυτών, τις αμφίδρομες μετατροπές συντεταγμένων, την παραγωγή πινάκων, τη φόρτωση των εγγενών πινάκων προτύπων γωνιών και ακμών, βασικές περιπτώσεις του εγγενούς βέλτιστου περιβλήματος και τον κύβο 2×2×2. Οι δοκιμές δεν επιχειρούν να αποδείξουν ότι κάθε κατάσταση 3×3×3 λύνεται γρήγορα τοπικά. Αντίθετα, ελέγχουν ότι κάθε υλοποιημένο τμήμα τηρεί το συμβόλαιό του και ότι η ετικέτα \code{exact} εμφανίζεται μόνο όταν η παραδεκτή αναζήτηση ολοκληρώνεται. Η διατύπωση αυτή είναι σημαντική: η επιτυχία των δοκιμών είναι αναγκαία, όχι ικανή, συνθήκη για την ορθότητα των ισχυρισμών της εργασίας.

Οι δοκιμές συντεταγμένων είναι ιδιαίτερα σημαντικές, επειδή οι πίνακες αποκοπής εξαρτώνται από ορθούς κωδικοποιητές. Ένα σφάλμα μετατόπισης κατά ένα στην κατάταξη θα μπορούσε να δημιουργήσει πίνακα που φαίνεται σωστού μεγέθους, αλλά δίνει λανθασμένα φράγματα. Οι δοκιμές συγκρίνουν τις μεταβάσεις των πινάκων κινήσεων με άμεσες κινήσεις κυβιδίων και ελέγχουν τις φασικές προβολές. Οι δοκιμές των επιλυτών ελέγχουν ρηχές περιπτώσεις, όπου η λύση επαληθεύεται γρήγορα. Οι δοκιμές του 2×2×2 ελέγχουν μικρή κατανομή περιορισμένου βάθους και ακριβείς αντιπροσωπευτικές λύσεις.

Η σχέση των δοκιμών με τους ισχυρισμούς του κειμένου είναι άμεση. Όταν το κείμενο λέει ότι υπάρχει ανεξάρτητος μηχανισμός επαλήθευσης, υπάρχει δοκιμή που ασκεί τη συμπεριφορά επαλήθευσης. Όταν λέει ότι υπάρχουν παραγόμενοι πίνακες αποκοπής, υπάρχουν δοκιμές και σενάρια που τους παράγουν. Όταν λέει ότι οι γωνιακοί και εξακμικοί πίνακες προτύπων του 3×3×3 είναι δυαδικά παραγόμενα τεκμήρια, υπάρχουν εγγενείς αμφίδρομες δοκιμές περιορισμένου βάθους και ο μηχανισμός επαλήθευσης ελέγχει τα πλήρη τεκμήρια. Όταν λέει ότι η μελέτη του 2×2×2 είναι πλήρης, ελέγχονται το πλήθος των καταστάσεων και η πληρότητα της κατανομής. Αν μια μελλοντική αλλαγή παραβιάσει αυτές τις ιδιότητες, οι αυτοματοποιημένοι έλεγχοι αποτυγχάνουν, ώστε η ασυνέπεια να εντοπίζεται πριν το κείμενο πάψει να αντιστοιχεί στα τεκμήρια.
